\newpage
\section{Conclusão}

\par Este trabalho partiu de um problema concreto. Os sistemas supervisórios convencionais exigem licenças proprietárias, equipamentos dedicados e instalação local especializada, o que os afasta de empresas de pequeno e médio porte, embora a necessidade de acompanhar, comandar e registrar um processo a distância seja a mesma de uma grande planta \cite{ZANGHI2019_SCADA_CONCEITOS,PHUYAL2020_COST_EFFICIENT_SCADA}. A resposta proposta foi um sistema supervisório web, multi-cliente e orientado a serviços, executado em contêineres, que o usuário acessa pelo navegador sem instalar nada. O objetivo geral, desenvolver e validar esse sistema, foi atendido na medida descrita nas seções seguintes: o sistema foi construído por inteiro, executa como um conjunto de contêineres em um único servidor e foi exercitado por testes automatizados, por campanhas de carga e por ensaios com equipamentos reais.

\subsection{Retomada dos objetivos}

\par O primeiro objetivo específico era revisar os conceitos que embasam a solução. O Capítulo~2 reuniu a fundamentação sobre sistemas supervisórios e seus módulos clássicos, a Indústria~4.0 e as tecnologias correlatas, a computação em nuvem, os contêineres e sua orquestração, os bancos de dados e o cache, as filas de trabalho, o protocolo Modbus em suas variantes \sigla{tcp} e \sigla{rtu}, a autenticação por tokens e a execução isolada de código. Essa base sustentou cada decisão tomada nos capítulos seguintes e permitiu ler o problema com o vocabulário da literatura \cite{COLOMBO2014_STATE_ART_AUTOMATION,LASI_INDUSTRY_40}.

\par O segundo objetivo era especificar os requisitos e projetar a arquitetura. O Capítulo~3 fixou dezenove requisitos funcionais e treze não funcionais e mostrou como cada módulo de um supervisório é atendido por componentes concretos: a interface web, a \sigla{api}, os \emph{workers} por família de fila, o serviço de aquisição, o motor de scripts isolado, o banco PostgreSQL, com um banco central e um banco próprio para cada cliente, e o Redis. A separação em serviços de implantação independente, ligados por mensagens e por filas duráveis, é a característica que torna o conjunto escalável e substituível por partes, e o fato de todos os serviços executarem em contêineres é o que o torna portátil entre servidores.

\par O terceiro objetivo era implementar a aquisição via Modbus TCP, com o agendamento das leituras e escritas e o registro do valor atual e do histórico. O Capítulo~4 descreveu o caminho do dado do campo ao banco e do banco à tela. O serviço de aquisição lê os dispositivos por afinidade, o valor corrente de cada tag fica no Redis, o histórico é gravado em lotes numa tabela particionada por mês e as médias por minuto e por hora atendem às consultas longas. No Capítulo~7, o ensaio de capacidade contra um servidor Modbus de teste sustentou até 9.750 tags a 30 s por réplica de aquisição, e a gravação do histórico chegou a cerca de 9.400 linhas por segundo para um cliente, com folga sobre a carga de um cliente típico. O isolamento entre clientes, com banco próprio e criação automática em poucos segundos, foi medido e atendido.

\par O quarto objetivo era a interface de supervisão. O Capítulo~5 apresentou as telas de cadastro, de acompanhamento em tempo real, o dashboard, o sinótico com editor gráfico, o mapa das plantas, os alarmes, a linha do tempo e o histórico com tendência, todas montadas sobre um conjunto único de componentes e guiadas pelos princípios de interfaces para automação \cite{ISA101_2015_HMI,ELIPSE_IHM_ALTA_PERFORMANCE}. O desempenho do lado do cliente, medido em celular, e o tempo de 8 ms para aplicar um ciclo de valores no sinótico indicam que a interface acompanha o processo sem sobrecarregar o navegador.

\par O quinto objetivo era o comando de equipamentos em quatro fases. O fluxo de seleção, verificação, confirmação e execução foi implementado com trava por tag, idempotência e registro de auditoria, e o sucesso do comando é definido pela resposta do controlador à escrita, sem releitura posterior do ponto. Os testes automatizados e a prova contra o servidor de teste exercitam toda a cadeia. A validação com o equipamento real, porém, depende dos ensaios de bancada descritos adiante, e por isso o requisito aparece como parcialmente atendido.

\par O sexto objetivo era o controle de acesso por cargo e por grupo. O sistema separa quatro cargos fixos e permite que cada cliente monte grupos de permissões a partir de um catálogo fechado, com a fronteira rígida no cliente. As senhas são guardadas com Argon2id, com custo de cerca de 24 ms por verificação, e os segredos que o sistema precisa ler de volta são cifrados. Esse conjunto foi coberto por testes de cargos, de grupos, de isolamento entre clientes e por um teste de segurança de caixa preta no fim de cada lote de desenvolvimento \cite{ARGON2_RFC9106,OWASP_CHEATSHEETS}.

\par O sétimo objetivo era validar o sistema em uma bancada com um \sigla{clp} e um inversor de frequência e medir o desempenho frente ao volume esperado. O Capítulo~6 descreveu a bancada, formada por um \sigla{clp} TM221 que simula o nível de um reservatório e por um inversor ATV320 que aciona um motor, e fixou o roteiro de seis ensaios. O Capítulo~7 apresentou o ambiente, as medidas de desempenho e a validação dos requisitos, e o conjunto dos resultados permite afirmar que a arquitetura cumpre o que se propôs na escala medida. Os ensaios de bancada propriamente ditos ainda dependem da execução pelo autor, como se explica a seguir.

\subsection{Contribuições}

\par A primeira contribuição é um sistema supervisório completo e funcional, de código próprio, que cobre o cadastro hierárquico, a aquisição, o histórico, os alarmes, os comandos e as telas de supervisão, e que pode ser executado com um conjunto de contêineres em um computador comum. A segunda é a organização multi-cliente com banco de dados próprio por cliente, criado e removido de forma automática, que oferece um isolamento mais forte que o de tabelas compartilhadas \cite{CHONG2006_MULTITENANT}. A terceira é o comando em quatro fases com garantias explícitas, cada uma associada a um mecanismo verificável por teste.

\par A quarta contribuição é a execução isolada do código escrito pelo usuário, em variáveis derivadas, eventos e componentes de conteúdo. O código roda em um interpretador compilado para WebAssembly, sem rede, sem arquivos e sem segredos, com limites de tempo e de memória \cite{HAAS2017_WEBASSEMBLY,WASMTIME_DOCS}. Isso dá flexibilidade ao integrador sem entregar a ele o servidor. A quinta é a documentação quantitativa do comportamento do sistema em recursos modestos, com tetos de capacidade por réplica e análise dos gargalos, que serve de base para dimensionar uma implantação e para comparações futuras.

\subsection{Limitações}

\par Os resultados devem ser lidos com as limitações que o próprio Capítulo~7 registrou. Todas as medidas foram obtidas em um único computador, que dividia os recursos com a carga e com outras atividades, e variaram até 50\,\% entre execuções, de modo que os valores são ordens de grandeza e não garantias. A carga de leitura veio de simuladores de resposta rápida, e o tempo de modems e de redes industriais reais reduz o teto prático. O efeito de réplicas adicionais não foi medido, e a redundância entre servidores é descrita e não implantada.

\par Entre os requisitos, sete aparecem como parcialmente atendidos. Os requisitos de comando em quatro fases e de sua conclusão pela resposta do equipamento (RF07 e RF09), de geração de alarmes (RF18) e de baixa latência das telas (RNF01) dependem dos ensaios com a bancada. Os de replicação horizontal (RNF02), de enfileiramento justo entre clientes (RNF03) e de disponibilidade (RNF04) dependem de medidas com mais de uma réplica e de ensaio de queda de réplica, que não foram feitas.

\par A capacidade do contêiner de \emph{workers}, que ficou abaixo do piso planejado para variáveis derivadas e tags internas, e a sensibilidade do isolamento entre clientes na fila de histórico à disputa de processador também permanecem como pontos a melhorar. Por fim, parte da avaliação de qualidade apoia-se em contagens de arquivos de teste e não em cobertura medida de ponta a ponta.

\todo[inline]{Fechar esta seção depois dos ensaios E1 a E6 da bancada: informar quais requisitos passam de parcial para atendido, os tempos de resposta do comando e o atraso do supervisório medidos, e ajustar a frase sobre o objetivo de validação em bancada. Se o Capítulo 7 for atualizado, conferir também a lista de requisitos parciais.}

\subsection{Trabalhos futuros}

\par O primeiro caminho de evolução é levar o sistema a um provedor de nuvem e a um orquestrador de contêineres. Como todos os serviços já executam em contêineres, a mudança não exige alterar o código, e o orquestrador oferece distribuição automática de réplicas, recuperação de falhas e atualização gradual \cite{NETTO2017_COORDENACAO_CONTAINERS_KUBERNETES,KUBERNETES_DOCS}. Essa implantação permitiria medir o efeito do número de réplicas sobre a vazão e comparar a execução em contêineres com a execução direta, a mesma comparação feita pelos trabalhos do orientador sobre contêineres gerenciados pelo Kubernetes \cite{NETTO2016_STATE_MACHINE_REPLICATION_KUBERNETES,NETTO2019_RAFT_KUBERNETES}.

\par O segundo caminho é a replicação do banco de dados e dos serviços de estado por meio de replicação de máquinas de estados e de um algoritmo de consenso como o Raft, o que tornaria o sistema tolerante à queda do servidor que hoje guarda o banco e o Redis \cite{ONGARO2014_RAFT,NETTO2019_RAFT_KUBERNETES}. A avaliação desse ganho, com injeção de falhas e medida da latência sob falha, seguiria o método das referências citadas.

\par Do lado dos protocolos, a arquitetura do serviço de aquisição prevê a inclusão de novas famílias de fila, o que abre espaço para o OPC~UA e o MQTT ao lado do Modbus TCP \cite{CARBALLO2024_MODERN_SCADA_OPCUA}. Do lado das funções, ficam como evolução a emissão de relatórios, as notificações por canais externos como o correio eletrônico e mensageiros, a restrição de acesso por planta, que hoje não existe porque a fronteira de acesso é o cliente, e um aplicativo móvel que explore a notificação ativa de alarmes. Também fica como trabalho futuro a execução dos ensaios de bancada com cargas reais mais variadas, inclusive com mais de um equipamento e com redes industriais de campo, para ampliar a confiança nos tetos medidos.

\par Em síntese, o trabalho entrega uma base funcional e verificável para um supervisório web acessível a empresas que não podem arcar com a solução tradicional. As limitações apontadas definem o próximo ciclo de trabalho, e a arquitetura em contêineres, escolhida desde o início, é o que permite que esse ciclo ocorra sem reescrever o que já foi construído.
